\documentclass[aps,prb,onecolumn,nofootinbib,superscriptaddress]{revtex4-2}
\usepackage{amsmath,amssymb,amsthm,mathtools,bm}
\usepackage{booktabs}
\usepackage[colorlinks=true,linkcolor=blue,citecolor=blue,urlcolor=blue]{hyperref}
\usepackage{microtype}
\newcommand{\Id}{\mathbf{1}}
\newcommand{\Ch}{\operatorname{Ch}}
\newcommand{\tr}{\operatorname{tr}}
\newcommand{\op}{\mathrm{op}}
\newcommand{\br}{\boldsymbol r}
\newcommand{\bk}{\boldsymbol k}
\newcommand{\ket}[1]{\lvert#1\rangle}
\newcommand{\bra}[1]{\langle#1\rvert}
\newcommand{\kett}[1]{\lvert#1\rangle\!\rangle}
\begin{document}
\title{Bulk topology from truncated overcomplete Wannier modes}
\author{Asadullah Bhuiyan}
\noaffiliation
\date{September 29, 2026}
\begin{abstract}
We construct a half-filled reference state from truncated overcomplete Wannier
modes and give a sufficient, volume-independent condition for its Chern number
to remain one. We then distinguish this static statement from a possible
deformation between imaginary-time filtering and monitored trajectories.
\end{abstract}
\maketitle
\setcounter{tocdepth}{2}
\tableofcontents

\section{The untruncated projector}
We consider the translation-invariant bulk, with two orbitals per unit cell,
and the parent Hamiltonian of Ref.~\cite{BPJ}. Its lower-band projector is
\begin{equation}
 P_-(\bk)=\tfrac12[\Id-\widehat{\boldsymbol n}(\bk)\cdot\boldsymbol\sigma],
 \qquad \boldsymbol n=(\sin k_x,\sin k_y,\alpha-\cos k_x-\cos k_y),
 \quad 0<\alpha<2 . \label{eq:parent}
\end{equation}
Here $\widehat{\boldsymbol n}=\boldsymbol n/|\boldsymbol n|$ is a unit vector,
not an operator. The half-filled ground state occupies $P_-$, with
$\Ch(P_-)=1$ in our orientation convention. Choose
$\tau_{A/B}=(1,\pm1)^{\mathsf T}/\sqrt2$ and write
$\ket{w_{\br\nu}}=P_-\ket{\br,\tau_\nu}$.
Completeness of the local trial states gives
\begin{equation}
 \sum_{\br,\nu}\ket{w_{\br\nu}}\bra{w_{\br\nu}}=P_-^2=P_-.
 \qquad
 \|w_{\br\nu}\|^2=\int_{\mathrm{BZ}}\frac{d^2k}{(2\pi)^2}
 \tau_\nu^\dagger P_-(\bk)\tau_\nu=\tfrac12 . \label{eq:tight}
\end{equation}
The last equality follows from oddness of $n_x/|\boldsymbol n|$.
Thus the normalized overcomplete Wannier (OW) columns
$W_{\br\nu}=\sqrt2\,w_{\br\nu}$ satisfy
$P_-=\frac12\sum_{\br,\nu}W_{\br\nu}W_{\br\nu}^\dagger$.
The factor $1/2$ is essential: this is a tight frame, not an orthonormal basis.
In the report's two-point ordering,
$C_{ij}=\operatorname{Tr}(\hat\rho\hat c_i^\dagger\hat c_j)$ and
$\Gamma=C^{\mathsf T}$; the ground state has $\Gamma=P_-$.

\section{Truncation and quantization}
Let $g_\ell(\br)$ be the indicator of $|x|,|y|\le\ell$, where
$\ell=n_{\rm shell}$, and normalize
$W^{(\ell)}_{\br\nu}=g_{\ell,\br}W_{\br\nu}/\|g_{\ell,\br}W_{\br\nu}\|$.
Introduce the synthesis matrices
$S_0=(W_{\br\nu})/\sqrt2$, $S_\ell=(W^{(\ell)}_{\br\nu})/\sqrt2$.
Since $S_0S_0^\dagger=P_-$, $\|S_0\|_\op=1$. Setting
$E_\ell=S_\ell-S_0$, $\epsilon_\ell=\|E_\ell\|_\op$, we obtain
\begin{align}
 F_\ell&=S_\ell S_\ell^\dagger,
 &F_\ell-P_-&=S_0E_\ell^\dagger+E_\ell S_0^\dagger+E_\ell E_\ell^\dagger,
 \nonumber\\[-2pt]
 \|F_\ell-P_-\|_\op&\le\delta_\ell,
 &\delta_\ell&=2\epsilon_\ell+\epsilon_\ell^2. \label{eq:bound}
\end{align}
Small column errors alone would not give a volume-independent bound.
Here translation symmetry does: write $W_\nu=W_{\boldsymbol0\nu}$,
$M=\max_\nu\|W_\nu\|_1$ and
$a_p=\max_\nu\|(1-g_\ell)W_\nu\|_p$, where the norms include both orbitals.
For $a_2<1$, normalization and the Fourier triangle inequality imply
\begin{equation}
 \epsilon_\ell=\sup_{\bk}\|E_\ell(\bk)\|_\op
 \le b_\ell:=a_1+
 \left[(1-a_2^2)^{-1/2}-1\right]M\longrightarrow0 . \label{eq:tail}
\end{equation}
Indeed, each normalized-column difference has $\ell^1$ norm at most
$b_\ell$; the two Bloch columns, each divided by $\sqrt2$, have combined
Frobenius norm at most $b_\ell$. Exponential OW localization makes both tails
vanish as $\ell\to\infty$, while $M$ is finite.

Assume $\delta_\ell<1/2$. Along $F(u)=P_-+u(F_\ell-P_-)$, spectral stability gives
\begin{equation}
 \operatorname{spec}F(u)\subset[-u\delta_\ell,u\delta_\ell]
 \cup[1-u\delta_\ell,1+u\delta_\ell],\qquad
 \Pi(u)=\mathbf1_{(1/2,\infty)}(F(u)),\quad 0\le u\le1. \label{eq:path}
\end{equation}
The occupation gap stays open. Equivalently,
$\Pi(u,\bk)=(2\pi i)^{-1}\oint_{|z-1|=1/2}(z\Id-F(u,\bk))^{-1}dz$,
with positive contour orientation. This defines a smooth, periodic rank-one
projector. Its invariant is
\begin{equation}
 \Ch(\Pi(u))=\frac1{2\pi i}\int_{\mathrm{BZ}}d^2k\,
 \tr\!\left(\Pi(u)[\partial_{k_x}\Pi(u),\partial_{k_y}\Pi(u)]\right).
 \label{eq:chern}
\end{equation}
On overlapping momentum patches, $\ket{v_a}=e^{i\theta_{ab}}\ket{v_b}$ and
$\bra{v_a}d v_a\rangle=\bra{v_b}d v_b\rangle+i\,d\theta_{ab}$.
Stokes' theorem identifies Eq.~\eqref{eq:chern} with the integer winding
$(2\pi)^{-1}\oint d\theta_{ab}$, summed over oriented patch boundaries.
It is also continuous in $u$, hence constant:
\begin{equation}
 \boxed{\epsilon_\ell<\sqrt{3/2}-1\quad\Longrightarrow\quad
 \Ch(\Pi_\ell)=\Ch(P_-)=1,\qquad \Pi_\ell=\Pi(1).} \label{eq:claim}
\end{equation}
Uniform truncation preserves momentum diagonality. The real-space projector
is exponentially local because it is the gapped spectral projection of a
finite-range operator. Its infinite-bulk three-sector invariant is the same
integer~\cite{Kitaev}; a finite disk still has finite-radius corrections.
This proves the claim for sufficiently large finite $\ell$, not specifically
for $\ell=1$. Nor is $F_\ell$ itself a projector or an occupation matrix:
its eigenvalues can exceed one. The reference Hamiltonian
$h_\ell=\Id-2F_\ell$ has half-filled ground-state projector $\Pi_\ell$.

\section{Imaginary time and the monitored circuit}
For a finite system, normalized $e^{-\beta\hat H_\ell}\kett{\psi_0}$ approaches
this ground state as $\beta\to\infty$, provided the initial half-filled state
has nonzero ground-state overlap; here
$\hat H_\ell=\sum_{ij}(h_\ell)_{ij}\hat c_i^\dagger\hat c_j$.
There is a useful connection to occupation measurements. With
$\hat N_{j,\pm}=\hat\chi_{j,\pm}^\dagger\hat\chi_{j,\pm}$ for normalized truncated modes,
the penalties $\hat q_{j,-}=\hat{\mathbf{1}}-\hat N_{j,-}$ and
$\hat q_{j,+}=\hat N_{j,+}$ are projectors, so
\begin{equation}
 e^{-\lambda\hat q}=\hat{\mathbf{1}}+(e^{-\lambda}-1)\hat q,
 \qquad e^{-\lambda\hat q_{j,-}}\xrightarrow{\lambda\to\infty}\hat N_{j,-},
 \qquad e^{-\lambda\hat q_{j,+}}\xrightarrow{\lambda\to\infty}
 \hat{\mathbf{1}}-\hat N_{j,+}. \label{eq:filter}
\end{equation}
These are precisely the desired-outcome, no-correction branches. Trotterizing
the positive penalty $\hat H_{\rm pen}=\sum_{j,\sigma}\hat q_{j,\sigma}$
gives $e^{-\beta\hat H_{\rm pen}}=\lim_{M\to\infty}
(\prod_{j,\sigma}e^{-\beta\hat q_{j,\sigma}/M})^M$ at fixed finite volume
and $\beta$. Its one-body matrix is
$2(F_+^{(\ell)}-F_-^{(\ell)})$, not generally $h_\ell$.
Here $F_-^{(\ell)}=F_\ell$ and $F_+^{(\ell)}$ is the analogous upper-band frame
operator. If $\|F_\pm^{(\ell)}-P_\pm\|_\op\le\delta_\pm$ and
$\delta_++\delta_-<1$, its ground-state projector is likewise connected to
$P_-$ without closing the gap. This is a static spectral argument, distinct
from dissipative preparation~\cite{Goldstein}.

We may therefore conjecture that typical late-time trajectories remain in
this deformed topological sector. However, weak-filter Trotterization does
not justify the strong-filter limit at fixed step size: the penalties do
not commute. An arbitrary record also contains the correction branches
$\hat\chi_{j,-}^\dagger$ and $\hat\chi_{j,+}$, which change physical charge;
projective branches are singular. Imaginary time under $\hat H_\ell$ preserves
charge. A deformation to an arbitrary record is consequently not automatic.
It would require control of bulk locality and a spectral or mobility gap
along a record-dependent interpolation, with any charge-changing localized
defects treated explicitly. Equation~\eqref{eq:claim} proves topology of the
reference state; extending it to the conditional ensemble is the open step.

\begin{thebibliography}{3}
\bibitem{BPJ} A. Bhuiyan, H. Pan, and C.-M. Jian,
\href{https://doi.org/10.1103/nk38-ygyq}{Phys. Rev. Research \textbf{8}, 023147 (2026)}.
\bibitem{Kitaev} A. Kitaev,
\href{https://doi.org/10.1016/j.aop.2005.10.005}{Ann. Phys. \textbf{321}, 2 (2006)}, Appendix C.3.
\bibitem{Goldstein} M. Goldstein,
\href{https://doi.org/10.21468/SciPostPhys.7.5.067}{SciPost Phys. \textbf{7}, 067 (2019)}.
\end{thebibliography}
\end{document}
